\documentclass[a4paper,11pt,openany]{ltjsbook}
\usepackage{luatexja}
\usepackage{luatexja-fontspec}
\usepackage{amsmath,amssymb,amsfonts,amsthm}
\usepackage{geometry}
\geometry{top=24mm,bottom=24mm,left=24mm,right=24mm,footskip=10mm}
\usepackage{tikz}
\usepackage{tcolorbox}
\tcbuselibrary{skins,breakable}
\usepackage{hyperref}
\usepackage{fancyhdr}

\definecolor{primaryblue}{RGB}{20, 60, 110}
\definecolor{softblue}{RGB}{238, 244, 252}
\definecolor{darkslate}{RGB}{35, 45, 60}

\hypersetup{
    colorlinks=true,
    linkcolor=primaryblue,
    citecolor=darkslate,
    pdftitle={古典で学び直すホワイトヘッド数学入門},
    pdfauthor={教育情報分析研究会}
}

\pagestyle{fancy}
\fancyhf{}
\fancyhead[LE,RO]{\thepage}
\fancyhead[RE]{\leftmark}
\fancyhead[LO]{\rightmark}
\setlength{\headheight}{14pt}
\renewcommand{\headrulewidth}{0.4pt}

\title{
    {\huge\bfseries\gtfamily 古典で学び直すホワイトヘッド数学入門}\\[14pt]
    {\large 底本：\textit{An Introduction to Mathematics} (Project Gutenberg eBook \#41568)}
}
\author{\textbf{教育情報分析研究会} 編著}
\date{2026年10月}

\begin{document}
\maketitle
\frontmatter
\tableofcontents
\clearpage
\mainmatter

\chapter{Chapter I: The Abstract Nature of Mathematics}

\#\#\# Chapter I: 数学の抽象的な性質

数学の研究はしばしば失望から始まる傾向がある。科学が重要な応用、概念的な興味深いアイデア、そして論理的な厳密さのある手法により、迅速に興味深い過程への導入を期待させる。しかし、ハムレットの父の幽霊のように、この偉大な科学は私たちの精神的な武器による取り組みを逃れてしまう——「ここにいる、そこにもいる、そして消えた」——そして見えてくるものは、幽霊が粗雑な方法で説明できないほどの優れた存在であるという理由では説明されないような虚構性を示す。

この科学が期待通りにならない理由は、その基本的な概念が学生に技術的な手続きから解き放たれて説明されていないためである。したがって、不幸な学習者は、全体的な概念が光っていない多くの詳細を理解しようと必死になる。もちろん、技術的な能力は価値ある精神活動の第一条件である：ミルトンのリズムやシェリーの情熱を理解するまでに、単語を拼写し、個々の文字の形が確実にある間、その詩の美しさを十分に評価することはできない。この意味では、学習には王道は存在しない。しかし、技術的な過程にのみ注目し、全体的な概念を考えることから除外することも誤りである。ここに、教条主義への道がある。

\#\#\# 本章の目的

次の各章の目的は、数学を教えることではなく、学生がコースの最初から科学が何であるか、そしてなぜ自然現象に対する正確な思考の基礎となるべきかを理解できるようにすることです。その後で述べられる科学内の詳細な推論への参照は例示のためのみであり、注意深く作られた一般的な議論は、読者が理解できない技術的な過程や記号が挙げられても、その内容が理解できるよう配慮されます。

多くの人々が数学と最初に出会う方法は算術を通じてです。二加二が四であることは通常、誰もが耳にしたことがある単純な数学者の命題の型例として取り上げられます。そのため、算術を考慮することで、可能であれば科学の最も明確な特性を見つけることができます。

現在、算術における最初に注目すべき事実は、その適用範囲が広いことです。味や音、リンゴや天使、心のアイデアや体の骨まで何でも適用されます。物事がどのようなものであれ、二加二が四であることには変わりありません。したがって、数学の主要な特性として、物事そのものが理由ではなく、何かが何かであるという性質と概念を扱うことを記述します。そしてこれらの性質やアイデアは、何かが何かであるという点から適用され、それらと直接関連のある感情や感覚とは無関係です。これが数学を抽象的な科学と呼ぶ意味です。

\#\#\# 抽象的な科学の重要性

抽象的な科学が人間の生活のことでほとんど重要でないと思われるのは自然なことだが、それは実際の興味深い要素を考慮に入れていなかったからである。スウィフトは「ゴルビル物語」においてラプタ王国の数学者の扱いについて二つの見解を持っている。彼はその国の数学家たちを馬鹿で無用な夢想家とし、彼らの注意を覚醒させるために拍子木が必要だと描写している。また、数学的な裁縫師は自分の身長を四分法で測り、他の寸法は規則とコンパスを使って導き出し、非常に不適切な服を作っている。

一方で、ラプタ王国の数学者は空中に浮かぶ磁気島という驚異的な発明によって国を支配し、その支配力を維持している。スウィフトはニュートンの「原理」がまだ書かれている時代には、現代世界を大きく変革する力となった一方で、数学家たちを嘲るような時とは特別に適していなかった。スウィフトが地震を笑うように、彼もまた当時の数学家たちを嘲笑してもよかったかもしれない。

しかし、数学の成果の一覧だけではその重要性を理解するには不十分である。数学がなぜ、その抽象性ゆえに常に最も重要な思考のテーマの一つでなければならないのか、少し考えることでその根本的な理由を明確にする価値がある。我々はなぜ事件の順序を説明する際には必ずしも数学的になる傾向があるのかを理解しようと試みる。

すべての出来事は相互に結びついており、それを考慮する必要があります。雷を観察したときには、それに伴う thunder を聞くことになります；風が吹くと海面に波が立つことを観察します；秋の寒さを感じると葉が落ちます。秩序はあらゆる所で支配しており、特定の状況が見つかった場合、他の状況も同様に存在すると予測できます。科学の進歩はこれらの相互関係を観察し、この動的な世界における出来事が少数の一般的な関係や法則の例であることを示すことにあります。個別の事象の中に普遍性を見出し、瞬時なものの中に恒常性を見出すことが科学的思考の目的です。科学の眼から見れば、リンゴが落ちる現象、太陽繞る惑星の運動、大気の地球への密着は全て重力法則の例として観察されます。最も複雑な瞬時的な状況を様々な恒常的な法則の例に分解する可能性は、現代思考の支配的アイデアです。

\#\#\# 法則の種類について

では、この科学的理想を完全に実現するために求められる法則の種類について考えてみましょう。私たちは周囲の世界の特定の事実に関する知識を得るために、感覚から情報を得ます。私たちは見たり、聞いたり、味わったり、嗅いだり、熱を感じたり、押し付けられたり、痛みを感じたり、しびれを感じたりします。これらはすべて私個人の感覚です：私の歯痛はあなたの歯痛ではありませんし、私の視覚もあなたの視覚でもありません。しかし、これらの感覚の起源を外部世界の事物間の関係に帰属させます。つまり、歯医者は歯痛を取り除くのではなく、歯そのものを取り除きます。さらに、私たちはすべての人々の観察が基盤となる一つの連続的な世界があることを想像しようと試みます。私の感覚とあなたの感覚のために異なる世界があるわけではありません；私たち両者が存在する一つの世界があります。患者と歯医者にとって同じ歯です。また、私たちは視覚、触覚、そして聴覚で同じ世界を観察します。

従って、これらの外部の事物との関係を特定の感覚に依存しない方法で説明したいということが容易に理解できる。外部の事物の世界における事象の進行が満足する法則は、盲人や耳聋の人々にとっても、私たちの理解を超えた能力を持つ存在にとっても、同じような中立的かつ普遍的な形で記述されるべきである。

しかし、直感的な感覚を捨て去ったとき、残されたものの中で最も有用な部分は、事物の抽象的な形式的な性質に関する一般的な概念から成り立っている。実際には、先に述べた抽象的な数学的概念である。従って、人間は段階的に進み、その過程の完全な意味を理解せずに、宇宙の性質に対する数学者的な説明を求めることになる。なぜなら、この方法以外では、特定の人々や特定の感覚タイプに依存しない事象の進行の大まかな概念を形成することは不可能だからである。

例えば、食事中に問われるかもしれない質問は：
「私のかんじ、あなたの触れる感じ、そして彼のかんじと臭いは何が基盤になっていたのか？」
答えは「リンゴ」である。

しかし最終的には、科学はリンゴを分子の位置と運動から説明しようとし、その説明は私やあなたや彼、視覚や触覚や味覚

\chapter{Chapter II: Variables}

\#\#\# Chapter II: Variables

数学は、古代ギリシャ人が数やいくつかの数についての命題を証明したときに科学として始まった。彼らが幾何学のために最初に述べた命題から、幾何学は偉大なギリシャ数学の一部となりました。しかし、代数は実質的な開始すらしなかった数世紀が経過しました。それでも、後のギリシャ数学家中にはいくつかの予兆がありました。

\#\#\#\# 1. 变量の導入

代数において、任意の数やある数を表すために文字を使用します。具体的な算術の数字ではなく、文字 \( x \) と \( y \) を使って、例えば \( 2 + 3 = 3 + 2 \) の代わりに、任意の二つの数 \( x \) と \( y \) について \( x + y = y + x \) という一般化をします。

また、\( 3 > 2 \) の代わりに、ある数 \( x \) についてある数やいくつかの数 \( y \) が存在し、\( y > x \) であるという一般化を行います。この後者の仮定は厳密な最終形で述べると仮定となるものですが、哲学と数学にとって極めて重要です。これは無限の概念を導入するためのものです。アラビア数字の導入により、文字が具体的な数を表すことを完全に廃止したことで、数学家たちに任意の数やある数という概念のために文字を使用することの技術的な利便性が示唆されました。

ローマ人はこの文書が記された年を MDCCCCX (1910) と表していましたが、我々は 1910 と書き、文字は他の用途に使用します。しかしこれは単なる推測です。

\#\#\#\# 2. 方程式の発展

代数の発展とともにニュートンとライプニッツによって微分 calculus が発明されましたが、これらの概念に関しては数学的思考の哲学における進歩は一時的に止まりました。そして最近の数年間で初めて、任意の数やある数が数学そのものに極めて根本的な役割を果たすことが認識されました。これによりさらなる数学的研究分野が開拓されることになりました。

\#\#\#\# 3. 基本的な代数的声明

以下に簡単な代数的声明をいくつか述べて、これらの基本的な概念がどのように現れるかを理解するために考えてみましょう。

\[
(1) \quad x + 2 = 2 + x
\]
\[
(2) \quad x + 2 = 3
\]
\[
(3) \quad x + 2 > 3
\]

最初に注意すべきは、「some」という言葉の含意にある可能性です。ここでの「any」はすべての数に対して成り立つことを示していますが、\( x + 2 = 2 + x \) は任意の数 \( x \) に対して真です。したがって、ある数 \( x \)についても真となります。「any」という言葉は「some」を含むことから、「some」は「any」を排しないことがわかります。

また、第二の例では実際には一つだけの数 \( x \) が存在し、\( x + 2 = 3 \) を満たすもので、それは1のみです。したがって、「some」はただ一つの数を指すこともあります。しかし第三の例では、任意の数 \( x \) が1より大きいとき、\( x + 2 > 3 \) が成り立ちます。これにより、この場合の「some」に対応する無限多種の数がわかります。したがって、「some」は「any」と「one only」の間で何でもあり、両端の境界ケースも含めることができます。

次に自然な形で式（2）と式（3）を質問に置き換えることができます：

式（2′） \( x + 2 = 3 \) を満たす数 \( x \)は何か？

式（3′） \( x + 2 > 3 \) を満たす数 \( x \)は何か？

考慮する式（2'）より、\( x + 2 = 3 \)は方程式であり、その解は容易に \( x = 3 - 2 = 1 \)であることが分かります。方程式 \( x + 2 = 3 \) の文脈から推測される問いに対する答えにおいて、\( x \)は未知数と呼ばれます。方程式の解き方は未知数を特定することです。

方程式は数学において非常に重要であり、式（2'）は元々の式（2）よりもはるかに根本的な概念を示しているように思われますが、これは完全な誤りです。不定な「変数」という概念が「何か」や「どれか」を使用する際に現れるものこそが数学において真に重要な概念であり、方程式における未知数は速やかに解くべきものとして重要ですが、その役割は補佐的なものです。

\#\#\#\# 4. 多変数の公式

特に「何か」の概念が含まれる場合には、多くの興味深い公式は一つ以上の変数を含みます。例えば、\( x + y = 1 \) を満たす分数または整数のペア \( x \) と \( y \) の考慮は、二つの関連する変数 \( x \) と \( y \) の概念を含んでいます。

二つの変数が存在する場合、主に二種類の声明が現れます。例えば、任意の数のペア \( x \) と \( y \) に対して \( x + y = y + x \)（式（1））、あるいはいくつかの数のペア \( x \) と \( y \) に対して \( x + y = 1 \)（式（2））です。

第二の種類の命題は、固定された関係によって結び付けられた数のペアの全体を考慮することを招く。この例では、\( x + y = 1 \)という関係によって結ばれた数のペアを考える。

最初のタイプの公式が任意の数のペアに対して真である場合、それらの公式は第二のタイプの公式を無限に多くの等価な形へと変形することができる。例えば、\( x + y = 1 \)という関係は以下の関係に等しい：
\[
y + x = 1, \quad (x - y) + 2y = 1, \quad 6x + 6y = 6,
\]
などである。したがって、命題を考慮している関係の適切な等価形は、その直近の目的にとって最も便利なものを使い、巧みな数学者はそれを用いる。

一般には、あるペアの項が固定された関係を満たす場合でも、一方の項が与えられたからといってもう一方の項が確定的に定められるわけではありません。例えば、\( y^2 = x \)という関係において \( x = 4 \) とすると、\( y \)は \( \pm 2 \)となり得ますので、任意の正の \( x \)に対して異なる値を持つ \( y \)が存在します。また、\( x + y > 1 \)という関係において、\( x \)や \( y \)のいずれかが与えられたからといってもう一方の項には無数の値が残ります。

\#\#\#\# 5. 関係式の領域

また重要な点が一つあります。正の数、整数または分数を考慮する場合、関係式 \( x + y = 1 \) を考えます。このとき、もし \( x \) や \( y \) のいずれかが \( 1 \) より大きい値を持つと、もう一方の数が正の数としてこの関係を満たす値は存在しません。したがって、関係式 \( x + y = 1 \) における \( x \) の「領域」は \( 1 \) より小さい数に制限され、同様に \( y \) の「領域」も制限されます。

\#\#\#\# 6. 整数の考慮

整数（正または負）のみを考慮し、関係式 \( y^2 = x \) を満たすような数のペアを考えます。このとき、

\[
y^2 = x
\]

という関係式から、例えば \( y = 3 \) の場合、\( x = 9 \) となります。逆に、\( x = 4 \) の場合、\( y \)は \( \pm 2 \) となり得ます。

これらの例を通じて、変数の概念と方程式の解き方がどのように数学を深く理解する助けになるかがわかります。これから先も、これらの基本的な概念を繰り返し練習することで、より深い洞察を得ることができるでしょう。

\chapter{Chapter III: Methods of Application}

\#\#\# Chapter III: Methods of Application

変数が関係を満たすという概念は、数学の応用において非常に重要な役割を果たします。この章では、具体的な例を通じてその重要性を深く理解しましょう。

\#\#\#\# 1. 建設費用の例

まず、建設費用に関する簡単な例から始めましょう。建設費用は1立方フィートあたり£1とします。さらに£20が1ポンドであるとします。新しい家を建てる複雑な状況の中で、所有者、建築家、builders（労働者）、見物人など、家が完成するまでのさまざまな感情や感覚の間、固定された関係は法によって保たれるものと仮定されます。

ここで、$x$ が立方フィートの数、£$y$ が費用であるとき、$20y = x$という関係が成り立つと仮定します。この$x$と$y$の関係は、どの所有者がどの家を建設する際にも真実であると仮定されます。

なぜこのような関係が成り立つのかを考えましょう。家を見たときにそれが何かを知らなければならないし、その建設に関連する出来事を認識しなければなりません。これらの出来事の中から、概念的に他の自然から隔離された場所で、コストと立方体の内容量の両方を決定することが可能になります。

例えば、$20y = x$という式が成り立つ場合、$x$ は立方フィートの数、$y$ はその立方フィートあたり£1の費用であることを表しています。つまり、家が100立方フィートの場合、建設費用は£200になります。

しかし、この法則は真実でしょうか？家を建設に関わった経験がある人なら、ここでのコストは高いと感じるでしょう。この価格で機能するのは高級なタイプの家だけです。これにより別の点が明確になる。公式 $20y = x$ に関連する数学的な計算をしている間、法則が真実であれ偽であれ、それは私たちにとって関係ありません。実際、$x$ と $y$ の意味付けは立方フィートの数とポンド Sterling の数であるとしても、それ自体には関係ありません。

数学的調査中では、実質的にこの変動する数字のペア $x$ と $y$ の間の相関性の性質を考慮しているに過ぎません。結果は、$y$ を漁師の数と解釈し、$x$ を捕れた魚の数と解釈しても同様に適用されます。つまり、平均して1人の漁師が20匹の魚を捕れるという法則が仮定されている場合でも、調査結果は公式 $20y = x$ の相関性の性質を与えるもので、数学的な確実性はその結果にのみ付与されます。

具体的な家を建設するコストには何の数学的確実性もありません。法則は完全には真実ではなく、それを提供する結果も正確とは限らない。実際、それは完全に間違っている可能性もあります。

\#\#\#\# 2. 牛顿と万有引力定律

次に、牛顿と万有引力定律について考えてみましょう。この法則は、任意の2つの物体が質量の積に比例し、それと逆比例する距離の平方に引き合う力をもつことを述べています。つまり、2つの物体の質量を$m$と$M$（lbs.で計測）とし、その間の距離を$d$（マイル）とするとき、一方の物体に対する他方からの引力は$\dfrac{mM}{d^{2}}$に比例します；したがって、この力を$k\dfrac{mM}{d^{2}}$と表すことができます。ここで$k$は、この引力の絶対的な大きさと、我々が選んだ力の尺度によって定まる一定の数です。

例えば、1lbs.の質量に対する重力として力を計測したい場合、$k$が代表する数値は非常に小さいことがわかります。なぜなら、$m$、$M$、および$d$をそれぞれ1に置き換えると、$\dfrac{kmM}{d^{2}}$は1lbs.の等しい質量間の距離1マイルでの引力となり、これはほとんど感じられません。

しかし、現在私たちは引力の力の公式を得ました。この力を$F$と呼びますとすれば、その式は$F = k\dfrac{mM}{d^{2}}$となります。これにより変数$F$、$m$、$M$、および$d$との関係が示されます。

ニュートンは orchard でリンゴが落ちるのを見ていて、その後、普遍的な重力法則が彼の頭に浮かび上がったとされています。最終的な法則の公式が彼の orchard で思いついただけでなく、他でも考えたかもしれません——彼にはどこかにいる必要があります。

ニュートンは、任意の質量間の力を表現する公式を考え出すために、膨大な前準備的思想や多くの頭脳と数世紀分の産物が必要でした。まず第一に、前の二章で説明した数学的な思考習慣や手続きが生成される必要がありました。それなしではニュートンは、任意の質量間の力を表現する公式を考え出すことはできなかったでしょう。

距離という最も簡単な用語から始めます。私たちは当然のようにすべての物質が一定の幾何学的な全体として考えられることを理解していますし、その各部分間の距離は特定の単位長さ（マイルやヤードなど）で測定できると考えています。これは私たちが最初に思いつく材料構造の最も基本的な側面です。

力と質量という他の用語は遥かに不明瞭です。ニュートンがこれらの言葉で伝えようとした思想の正確な理解はゆっくりと成長し、実際にはニュートン自身が初めて動力学の真の一般的原理を完全に把握した最初の人でした。

\#\#\#\# 3. 結論

これらを通じて、数学的な法則や公式は具体的な結果を提供しますが、その結果が自然の事実に正確に適用される確実性は必ずしもありません。議論の結論は、その始まりの仮定よりも決してより確実ではありません。

したがって、ある出来事が正確に予測されるほど正確な計算を行っても、その法則は真実であるかどうかの疑問は常に残ります。もし法則が具体的な結果を述べているなら、ほぼ確実にそれは正確ではありません；そして、最良の場合でも、厳密に計算された結果はおそらく起こりません。

しかし、我々には理想的な精度で観察する能力がないため、結局のところ、不正確な法則でも十分良いものであるかもしれません。

\chapter{Chapter IV: Dynamics}

\#\#\# Chapter IV: Dynamics

世界は十八世紀まで待たされ、ギリシャの数学物理学者がその後を引き継いだ。当時の16世紀と17世紀における偉大なイタリア人、特にアーリストオレス（画家）、レオナルド・ダ・ヴィンチ（生年1452、没年1519）とガリレオ（生年1564、没年1642）は、アーキメデスが知っていた秘密を再発見し、抽象的な数学的概念と自然現象の実験的調査との関連性を見つけることに成功した。一方で、数学の進歩と正確な天文学知識の蓄積により、自然哲学者は研究においてより有利な立場に置かれた。また、その時代の自我中心的な自己主張や個人経験への渇望は、思考者たちを自身で現象を見るために動かし、帰納的論理における数学理論と実験との関係の秘密が実質的に解明された。

ガリレオという哲学者がピサの傾斜塔から重りを落としたことは、その時代の特徴的な行為だった。常に思考者と行動者の存在がある。数学物理は、同じ人物の中に思考の衝動と行動の衝動を組み合わせた時代の産物である。

この塔から重さの異なる物体を落とすという事柄は、知識における重要な段階を象徴的に示しています。これは、全分野の基礎となる動力学に関する正しい概念を得た最初のステップに匹敵するものです。議論の焦点は、同じ高さから落とされた異なる重さの物体が同じ時間で落ちるかどうかでした。アリストテレスの定説によれば、その時代まで普遍的に受け入れられていたように、 heavier な物体の方が速く落ちるとされました。ガリレオはこれらが同じ時間で落ちることを主張し、彼の主張を立証するために、傾斜した塔の頂上から重さの異なる物体を落としました。

この規則に対する例外的なケースは、例えば極めて軽量や高速な理由により空気抵抗が重要となる場合に生じます。しかし、空気を無視すれば法則は正確です。

ガリレオの成功した実験は単なる運の良さによるものではなく、運動の不変性と質量に関する彼の正しい概念から生まれたものです。現在ニュートンが述べる最初の運動法則は次の通りです---全ての物体は、強制的な力によってその状態を変えるまで、静止状態または一直線の等速運動の状態を維持します。この法則は単なる乾いた公式以上のものであり、敗れた異端者に対する勝利の讃歌でもあります。法則から「または一直線の等速運動」の表現を取り除くことで、アリストテレスの反対法則として解釈できる形を得ることができます---全ての物体は強制的な力によってその状態を変えるまで静止状態を維持します。

この最後の誤った公式では、力以外に物体は静止状態を維持すると主張され、従って物体が動いている場合、運動を維持するためには力が必要であると述べられている。したがって、力を停止させると運動も停止する。真のニュートンの法則はこの点から全く反対の立場を取ります。力によって作用されない物体の状態は一直線の均一な運動であり、外部の力や影響を原因として探す必要はありません。または、あなたがそう定義したいなら、この均一な直線運動の不変の伴侶と見做すこともできます。静止はただそのような運動の一例に過ぎず、速度が常にゼロである場合だけです。したがって、物体が動いているときには、速度の率や方向の変化を説明するための外部の影響を探さないで済む場合があります。物体が同じ速度と同一の方向で動いている限り、力の援助を求めることはありません。

\#\#\# コペルニクスとケプラー

二つの視点の違いは、惑星の運動理論を通じてよくわかります。コペルニクスはポーランド人で、西普鲁士トーン出身（1473年生まれ、1543年に没）でした。彼は太陽を中心に近似円軌道上で回る惑星を含め地球も含めて考えると、それほどシンプルな説明になることを示しました。その後、ドイツの数学者ケプラーは1609年に実際にはこれらの軌道がほぼ楕円であることを証明し、これは後の章で詳しく説明する特別な種類の卵形曲線です。

すぐに問題となったのは、これらの力は何によって惑星がこのような運動を保たれているのかという点でした。旧来の誤った見解では、ケプラーは実際の速度そのものが必要な力をもって保存されるべきであると考えていました。したがって、添付図（図4）のように接線方向の力を求めました。しかしニュートンの法則によれば、一部の力以外は惑星は既存の速度で一直線に動くため、太陽から完全に離れてしまうことになります。そこでニュートンは、運動を楕円軌道へ曲げる力が必要であることを探求しました。彼は次のような図（図5）において、その力を太陽に向かう方向であると示しました。実際には、その力は距離の逆平方に従って働く太陽の引力です、これは先ほど述べた通りです。

ギリシャ人によって機械学の科学は、レバーを使用して得られる機械的利得に関する理論から発展し始めました。また、物体の重さに関連する様々な問題も考慮されました。16世紀末から17世紀にかけて、前段落で述べたように、この分野は主に落下する物体の理論を説明するためにも、そして主には惑星運動の科学的な理論を提供するために基礎づけられました。しかし、それ以来、動力学はより大きな任務を引き受け、他の分野とは異なる最終的な科学として自らを位置付け始めました。その主張は、物体の感覚で捉えられる様々な性質は、空間に存在するものの位置変化に対する我々の特殊な評価方法であるという点まで到達しています。

例えば、ウェストミンスター寺院を見ると、それは数世紀間灰色と動かないまま立っています。しかし、現代科学の理論によれば、その灰色が建物の不動性を強調するように感じさせるものは、実際には建築物の外表面を形成し、エーテルに振動を伝える最終的な分子の急速な運動の方法であるだけです。また、手で石を触ると、その冷たさと均一な温度を感じますが、これは建物の静かな安らぎを象徴しています。しかし、この感覚は手から石へ、または石から手への熱の移動に対する我々の認識を示すだけです。そして、現代科学によれば、熱は物体の分子の振動であるとされます。最後に、オルガンが演奏し始めると、音は耳の鼓膜に空気の運動が衝突して生じる結果です。

従って、現象を動力学的な説明で与える試みは、そのようなものや体がこの場所にあり現在は彼方の場所にあるという一般的形式の声明によってそれらを説明することである。これにより我々は現代科学の大切な基本的な概念へと到達する、すなわちすべての感覚は空間における事物の構造の変化との比較の結果である。従って運動の法則、つまり事物の構造の変化の法則は物理学の最終的な法則であることがわかる。

数学を自然哲学の研究に適用するとき、科学は系統的に普通の思考が偶発的に行うことを行います。椅子について話すと、通常は何らかの形で見たり感じたりしたものを指します；しかし私たちの言葉の大半は、視覚や感覚とは独立して何かが存在することを前提としています。一方で、数学的な物理学ではその反対の道筋を歩みます。椅子は特定の誰かや特別な感知モードに関連付けずに考えられます。結果として、思考の中で椅子は空間に位置する分子の集合や電子群、運動中の以太の一部などと見なされます。しかし重要なのは、科学が椅子を空間で動くものとして捉えることである。

\#\#\# 解説

この章では、動力学という概念について深く掘り下げていきます。まず、ガリレオの実験から始めてみましょう。ピサの傾斜塔からの落体実験は、当時の科学者たちにとって大きな挑戦でした。ガリレオは、重さの異なる物体が同じ時間で地面に落ちることを主張し、その証拠として実際に実験を行いました。

この実験は、アリストテレスの定説（ heavier な物体の方が速く落ちる）に対する挑戦でした。しかし、空気抵抗がある場合や極端な状況では、この法則が適用されないことがわかりました。ガリレオの成功は、運動の不変性と質量に関する彼の正しい概念から生まれたものです。

次に、コペルニクスとケプラーの貢献について触れましょう。コペルニクスは太陽中心説を提唱し、惑星が近似円軌道上を回ることを示しました。その後、ケプラーはこれらの軌道が楕円であることを証明しました。この発見は後の章で詳しく解説します。

ニュートンの運動法則についても触れています。彼の最初の法則によれば、全ての物体は強制的な力によってその状態を変えるまで、静止状態または一直線の等速運動の状態を維持します。これは単なる公式以上のものであり、アリストテレスの反対法則と対比しながら解説しています。

最後に、現代科学における動力学の概念について触れています。建物や椅子などの現象は、実際には分子レベルでの運動によって形成されていることを示唆しています。これは、我々が日常で感じている静止や温度を感じる方法とは異なる視点を提供します。

これらの解説を通じて、大人の読者が動力学の基本的な概念とその歴史的背景を理解し、新たな視点から自然現象を捉えることができるようになるでしょう。

\chapter{Chapter V: The Symbolism of Mathematics}

\#\#\# Chapter V: The Symbolism of Mathematics

\#\#\#\# 数学における記号の重要性

我々は、アラビック記法で $0$,~$1$, $2$,~\dots, $9$, $10$, $11$,~\dots\Add{,} $100$, $101$,~\dots\ と表される整数について十分な概念を持っていることを仮定しよう。この記法はアラブ人を通じてヨーロッパに導入されたが、彼らはおそらくヒンドゥーの源流からそれを得たものである。最初に系統的に説明された作品は、インドの数学家中の1人が著した「Bhaskara（1114年生）」によるものである。しかし実際の数字符号は私たちの紀元7世紀まで遡ることができるし、おそらく本来チベットで発明されたものかもしれない。しかし現在の我々にとって、記法の歴史は詳細な点に過ぎない。この数字システムが示す興味深い点は、良い表記法の巨大な重要性である。不要な作業から脳を解放することで、良い表記法はその力をより高度な問題に集中させ、実質的に人間の知能を増加させる効果がある。アラビック記法が導入される以前には乗算も難しかったし、整数の除算は最高の数学的才能を要求した。おそらく現代世界でギリシャの数学家が驚くべきことに、強制教育の影響下に西欧の大半の人々が大きな数に対する除算操作ができるようになった事実は、彼にとって絶対的な不可能性に等しいものだったろう。この記法を十進小数まで拡張したのは17世紀になってからである。我々現代の容易な計算力は、十進小数への適切な表記が徐々に発見される結果として、ほとんど不思議な結果である。

数学は多くの人々にとって難解で、神秘的な科学と見なされることが多い。その理由は、数多く用いられる記号によるものである。もちろん、理解できないシンボルはもっと不思議に思えるだろう。また、一部しか理解できず慣れていないシンボルは難しい。同じように、どの職業や仕事にも専門的な技術用語があり、それらは未訓練の者には難解である。しかし、これはその性質自体が難しいからではなく、むしろそれを簡単にするために導入されたものである。数学においても、真剣に数学的概念に取り組んでいる場合、シンボルは必ずしも簡略化するものではないが、それは大きな簡便さをもたらす。実用的なだけでなく、非常に興味深いものでもある。なぜなら、それは概念の分析であり、相互関係をほとんど図的に表現しているからである。もし誰かがシンボルの有用性に疑問を持った場合、以下の代数の基本法則をすべてのシンボルを使わず全文で書き出すことを試みてみることをお勧めする（これらは代数の基本法則を表す）：

\begin{gather*}
x + y = y + x, \Tag{(1)} \\
(x + y) + z = x + (y + z), \Tag{(2)} \\
x × y = y × x, \Tag{(3)} \\
(x × y) × z = x × (y × z), \Tag{(4)} \\
x × (y + z) = (x × y) + (x × z). \Tag{(5)}
\end{gather*}

ここで、式(\ref{1})と式(\ref{2})は加法の交換律と結合法則を表し、式(\ref{3})と式(\ref{4})は乗法の交換律と結合法則を表し、式(\ref{5})は加法と乗法の分配律を表している。例えば、シンボルを使わなければ、式(\ref{1})は「任意の数に二番目の数を足すことは、最初の与えられた数に二番目の数を足すことと同じ結果になる」という意味となる。

この例は、シンボルの助けによって論理を推移させることが、眼でほとんど機械的にできるということを示しています。それ以外の場合には、脳の高い機能が呼び出される必要があることを示しています。

「我々は自分たちが何をしているかを考える習慣を育てることを奨励する」というコピー本や有名な人々のスピーチでは繰り返される箴言は極めて誤ったものだと断言できます。しかし、その逆が正解です。文明は、考えずにできる重要な操作の数を増やすことで進歩します。

シンボルに非常に重要なものを持たせる一つの性質は、それが一目で理解でき、速やかに書き下せることです。シンボルをよりコンパクトに配置する方法としては、直接接続する以外にはありません。よって、良いシンボル表現において重要なシンボルの接続は重要な意味を持つべきです。これはアラビア数字による数値表記の一つの長所でもあります；十つのシンボル $0$,~$1$, $2$, $3$, $4$, $5$, $6$, $7$, $8$,~$9$ と単純な接続によって、どのような数も表現できます。また、代数学において、変動する二つの数 $x$ および $y$ の接続 $xy$ を何を表すかを選択しなければなりません。現在の手元にある最も重要なアイデアは加算と乗算です。数学的には、$xy$ を $x × y$ とするようにシンボル表現をコンパクトにすることを選びました。したがって、上記の法則 \Eq{(3)},~\Eq{(4)}, および \Eq{(5)} は一般的には以下のように書き表されます：
\[
xy = yx,\quad
(xy)z = x(yz),\quad
x(y + z) = xy + xz,
\]
これにより、大きなコンパクト性の利益を得ることができます。同じシンボル表現の規則は定数と変動する数の接続にも適用されます；$3x$ は $3 × x$, $30x$ は $30 × x$ のように書き表します。

明確に述べると、変数に具体的な数値を代入する際には、×の記号を適切に復元することが必要で、アラビア数字の表記と矛盾しないように注意しなければなりません。例えば、x に 2 を y に 3 を代入するとき、xy の代わりに 2 × 3 と書き、23 ではなくします。23 は 20 + 3 を意味します。

また、見かけ上は控えめな記号が科学の発展にとっていかに重要であるかを指摘するのは興味深いことです。そのような記号は重要なアイデアの強調表示を可能にし、しばしば非常に細かいアイデアを表現しやすくします。そしてそれが存在することで、そのアイデアと複雑な思想の流れとの関係性を容易に示すことができます。

例えば、最も控えめな記号である 0 を考えてみましょう。これは数零（ゼロ）を表します。古代世界の多くの数学家中のほとんどは、数零という概念が非常に困惑させたでしょう。なぜなら、それは非常に細かいアイデアで、明らかに直感的ではないからです。量の零についての意味に関する議論は哲学的な著作に多く見られます。実際には、ゼロは他の基数数と比べても、より難しく或いは細かい概念ではありません。1 または 2、または 3 とは何を意味するのか？これらは我々が日常的に使用しているアイデアですが、その基本的な要素について明確な分析ができる人は多くはありません。ゼロの重要な点は、日常生活での操作には必要とされないことです。誰も魚を買うときに零個を購入することはあり得ませんが、数学ではゼロは非常に重要です。

\#\#\#\# まとめ

この章で学んだように、数学における記号は単なる便利な道具ではなく、思考のフレームワークを形成する重要な要素です。良い表記法は計算や証明を容易にし、複雑な概念を簡潔に表現します。ゼロという概念も、その存在がなければ現代の数学は全く異なる形になっていたかもしれません。これから先の章では、これらの記号を使ってより高度な数学的概念を解説していきます。

この学びを通じて、あなた自身の思考プロセスを深く理解し、新たな視点から数学を見直す機会を得ることができます。ぜひ、この旅を楽しみながら進んでください。

\chapter{Chapter VI: Generalizations of Number}

\#\#\# Chapter VI: Generalizations of Number

数学の大きな特徴の一つは、整数から始まった概念群が発明されたことです。これらの概念は数を拡張または一般化したものと呼びましょう。まず最初に Fraction（分数）という概念があります。我々が所有する最も古い算術に関する論文は、紀元前1700年から紀元前1100年にかけてエジプトの司祭であるAhmesによって書かれました\index{Ahmes} おそらくもっと古い著作のコピーです。この論文は主に分数の性質について扱っています。したがって、この概念は数学の歴史において非常に早期に発展したことがわかります。実際、その主題は非常に自然なもので、田地を3等分し、そのうち2つを取り出すという操作はしばしば行われたでしょう。そのため、遠い文明の人々が三分の二という概念を熟知していたのは驚くことではありません。

そこで最初に数の一般化として分数の概念を置きます。ギリシャ人はこの主題を比という形で考えていたため、ギリシャ人にとって2フィートの線と3フィートの線との比は2対3であると言ったでしょう。我々の代数学記号の影響を受けた場合、一方の線が他方の三分の二であると言った方が自然であり、三分の二是数値的な乗数と考えることになります。

\#\#\# 1. 比と分数

ギリシャ人は比の理論や分数に関連して、大きな発見を行いました。この発見は哲学的な考えだけでなく、数学的な思考にも多くの影響を与えました。彼らは「無比値」の比の存在を発見しました。実際、彼らの幾何学的調査の中で、任意長さの線から他の線が存在することを証明しました。これらの線の長さは元の長さに対する整数ペアの比で表すことができません——言い換えれば、これらは元の長さの正確な分数ではない長さです。

例えば、正方形の対角線は同じ正方形の辺の任意の分数で表現することはできません。現代の記法では、対角線の長さは辺の長さの $\sqrt{2}$ 倍です。しかし、$\sqrt{2}$ という値を正確に表す分数はありません。$\sqrt{2}$ の値に近づくことはできますが、絶対的にその値には到達しません。例えば、$\dfrac{49}{25}$ は $2$ より少々小さい一方で、$\dfrac{9}{4}$ は $2$ より大きいので、$\sqrt{2}$ は $\dfrac{7}{5}$ と $\dfrac{3}{2}$ の間にあることがわかります。$\sqrt{2}$ の値に近づく最も systematic な方法として、小数分数の列を生成するには、平方根を通常の方法で求めるのが最良です。したがって、その列は $1$、$\dfrac{14}{10}$、$\dfrac{141}{100}$、$\dfrac{1414}{1000}$ などとなります。

このような比はギリシャ人によって「無比値」と呼ばれています。それらはギリシャ時代から現在まで、大きな哲学的な議論を引き起こし、その困難性は最近まで解消されていませんでした。

\#\#\# 2. 実数の概念

無比比例を分数と並べ、整数、分数、無比比例の全体を一つの数のクラスとして考えます。これを「実数」と呼びます。常に実数を大きさの順序で配置し、ゼロから上に進むにつれて無限大になるように考えます。実数は線上的な点で便利に表現できます。任意の線$OX$を選び、$O$から始まり、$OX$方向に無限に延びる線として定義します。その上で、適当な点$A$を選んで$OA$を1単位長さとします。次に、$AB$, $BC$, $CD$などの各辺をそれぞれ$OA$と同じ長さで分けていきます。このとき、点$O$は数$0$、$A$は数$1$、$B$は数$2$を表します。実際、任意の点が表す数はその点から$O$までの距離（$OA$を単位長さとして）の測定値です。$O$と$A$間の点は真分数や$1$未満の無比比例を、$OA$の中点は$\dfrac{1}{2}$、$AB$の中点は$\dfrac{3}{2}$、$BC$の中点は$\dfrac{5}{2}$を表します。このようにして$OX$上の任意の点が一つの実数を表し、逆に任意の実数も$OX$上の特定の一点で表現されます。

$OX$上から$O$を起点として、$O$から$X$に向かって規則的に配置された点の系列は、大きさの順序で昇順に並べられた実数を表します。

\#\#\# 3. 整数と分数との違い

すべてこれだけは単純に思えるが、この段階で興味深いアイデアを抽出することができる。整数を表す点の系列を考える。つまり、点$O$, $A$, $B$, $C$, $D$, 等である。ここでは最初の点$O$があり、次に続く明確な点$A$があるし、各点（例えば$A$や$B$）は一方には明確な直近の前駆点、他方には明確な直近の後継点を持つが、点$O$には前駆点がない。また、この系列は無限に延び、終了点がない。このような順序は整数の型の順序と呼ばれ、その本質は行の中の1番目の点以外の各点が両隣に近接した先駆点と後継点を持つことにある。

次に、整数と分数を一緒に考え、無理数に対応する点を省略する。このとき得られる系列の順序は全く異なるものである。最初の項$O$があるが、どの項も明確な直近の前駆点や後継点がない。これは容易に確認できる。例えば、任意の2つの分数の間に常に別の値の中間の分数を見つけることができる。その一つは分数を足し合わせてその結果を半分にする方法である。例えば、

\[
\frac{2}{3} \text{と}\frac{3}{4} の間に }\frac{17}{24} = \frac{1}{2} \left( \frac{2}{3} + \frac{3}{4} \right) \text{ が存在し、}
\]

\[
\frac{2}{3} \text{と}\frac{17}{24} の間に }\frac{33}{48} = \frac{1}{2} \left( \frac{2}{3} + \frac{17}{24} \right) \text{ が存在し、}
\]

このような無限に続く。この性質から、この系列は「密接」であると呼ばれる。系列には終点がないが、線$OX$を進めるにつれて無限に増大する。最初の見方では、分数から整数を含む系列とすべての実数（整数、分数、無理数を含む）から得られる系列は同じ型であるように思えるかもしれない。しかし、重要な違いがあり、それが今後開発される。

\#\#\# 4. 無比比例の重要性

分数の系列に無理数が含まれていないことで、特定のクラスには端点がない状態が生じる。例えば、無理数$\sqrt{2}$を考える。実数の系列では、これは全ての平方が$2$未満の数と$2$より大きい数を分ける位置にある。しかし、分数だけの系列に限定し、無理数を考慮しない場合、$\sqrt{2}$を導入できないので、どの分数もそのような端点を持つ性質を持たない。したがって、分数の系列には$\sqrt{2}$の位置に類似の欠陥がある。

このように、分数の系列に欠けている端点は小さな問題のように思えるかもしれませんが、数列の一部が全体の数列を覆わない場合でも極限や最大値の存在が保証されないことは大きな欠点であることを、どの数学家も知っている。

\chapter{Chapter VII: Imaginary Numbers}

\#\#\# Chapter VII: Imaginary Numbers

本章では、虚数という神秘的な数学的概念について深く掘り下げていきます。前章で扱った数学的概念も人気を博しましたが、虚数はその成功を全く異なる形で見せました。実用的な人々だけでなく、文学者や哲学者も、数学家が神秘的な存在であると認めながらも虚数に没頭することについて困惑していました。

\#\#\#\# 虚数の起源

虚数の概念は、正負の数と同じように、方程式を解くために生まれたものです。正負の数は方程式 \(x + 1 = 3\) や \(x + 3 = 1\) のような式から導かれました。同様に、虚数も方程式 \(x^2 + 1 = 3\) や \(x^2 + 3 = 1\) の解を求めるために生まれました。

例えば、方程式 \(x^2 + 1 = 3\) を解くと、
\[ x^2 = 2 \]
となり、この式の解は
\[ x = \pm\sqrt{2} \]
となります。これは問題ありませんが、方程式 \(x^2 + 3 = 1\) の場合、
\[ x^2 = -2 \]
となり、ここで困ったことに、正数や負数を自身で乗算しても負の結果は得られません。

\#\#\#\# 虚数の導入

そこで、方程式 \(x^2 + a = b\) の解を求めるために、新しい記号と解釈が必要になりました。特に、\(\sqrt{-1}\) という記号が導入されました。この記号は \(-1\) の平方根として定義され、虚数単位 \(i\) と呼ばれます。

\[ i = \sqrt{-1} \]

これにより、方程式 \(x^2 + a = b\) の解を
\[ x = \pm\sqrt{b - a} \]
のように表すことができます。ただし、\(a\) が負のとき、\(\sqrt{-c^2}\) と書き換えることで、
\[ \sqrt{-1} = i \]
となります。

\#\#\#\# 虚数の解釈

虚数 \(i\) の解釈は、正負の数や有理数、無理数と同じように形式的な法則を満たす必要があります。つまり、加法、減法、乗法、除法の通常の法則が成り立つはずです。

例えば、
\[ i^2 = -1 \]
という性質があります。さらに、虚数と有理数や無理数との操作も可能で、これらは一般的な代数法則を満足します。

\#\#\#\# 平面座標系での表現

虚数の概念を理解するためには、平面座標系での表現が役立ちます。直線 \(XOX'\) および \(YOY'\) を直角に交差させたものを「軸」として取り、すべての測定からここから始めるものとします。

正の数と負の数を利用することで、平面内の任意の点の位置は一対のこのような数で表現できます。例えば、順に書かれた一対の数 \((+3, +1)\) が存在し、この一対には最初の数（上方向）と二つ目の数（右方向）があります。

虚数 \(i\) を含む複素数は、平面座標系で点を表すのに使われます。例えば、複素数 \(3 + i\) は平面内の点 \((3, 1)\) に対応します。

\#\#\#\# 虚数の歴史的背景

虚数の概念が生まれたのは18世紀頃です。当時の数学家たちは、虚数を無意味な記号として扱っていましたが、後にこれらの記号に意味を与える解釈が有用であることが認識されました。

例えば、方程式 \(x^2 = -3\) は当初無意味とされていましたが、虚数 \(i\) を導入することで、
\[ x = \pm\sqrt{-3} = \pm i\sqrt{3} \]
という解を得ることができます。このように、虚数の概念は数学の発展に大きな役割を果たしました。

\#\#\#\# 結論

虚数は最初無意味とされていましたが、現在では重要な数学的概念として認識されています。虚数 \(i\) の導入により、方程式の解を完全に求めることが可能になりました。また、平面座標系での表現も虚数の理解を深めます。

虚数の概念は、正負の数や有理数、無理数と同じように形式的な法則を満たす必要があります。虚数 \(i\) の導入は数学の発展に大きな影響を与えました。これからは、虚数がどのようにして現代の数学に組み込まれたのか、そしてその重要性について深く理解しましょう。

この学び直しの旅を通じて、虚数という神秘的な概念も、実は私たちの日常生活や科学技術にも深く関わっていることを再認識することができます。

\chapter{Chapter VIII: Imaginary Numbers (Continued)}

\#\#\# Chapter VIII: Imaginary Numbers (Continued)

\#\#\#\# 1. 乗法の定義とその特性

まず、順序対の乗法について考えましょう。乗法は加法と同様に、いくつかの重要な特性を持つべきです。

- **(α) 結果は別の順序対となるべきです：** つまり、乗算の結果もまた順序対になるということです。
  
- **(β) 操作は可換性を持つので：**
  \[
  (x, y) × (x', y') = (x', y') × (x, y)
  \]
  これは、順序対の乗算が交換律を満たすことを意味します。

- **(γ) 操作は結合法則を満たし：**
  \[
  \{(x, y) × (x', y')\} × (u, v) = (x, y) × \{(x', y') × (u, v)\}
  \]
  これは、乗算が結合律を満たすことを意味します。

- **(δ) 除算の結果が一意的であることを保証する必要があります（ゼロ対 $(0, 0)$ の場合を除く）：** つまり、未知の順序対 $(x, y)$ を求め、方程式
  \[
  (x, y) × (a, b) = (c, d)
  \]
  を満たすとき、一意的な答えが存在し、これを次のように表すことができます：
  \[
  (x, y) = (c, d) ÷ (a, b), \quad \text{または} \quad
  (x, y) = \frac{(c, d)}{(a, b)}
  \]

- **(ε) さらに、加法と乗法を含む法則、すなわち分配律が満たされるべきです：**
  \[
  (x,y) × \{(a, b) + (c, d)\} = \{(x, y) × (a, b)\} + \{(x, y) × (c, d)\}.
  \]

これらの条件は、最初に複雑に見えるかもしれませんが、単純な幾何学的解釈を持つ定義によって満たすことができます。

\#\#\#\# 2. 乗法の具体的定義

順序対 $(x, y)$ と $(x', y')$ の乗算を以下のように定義します：
\[
(x, y) × (x', y') = \{(xx' - yy'), (xy' + x'y)\}
\]

この定義から明らかに、乗算の結果は別の順序対となることがわかります。また、方程式 $(A)$ の右辺の値は、同時に $x$ と $x'$、および $y$ と $y'$ を入れ替えることで変化しないこともわかります。したがって条件 (α) と (β) は明らかに満たされます。

\#\#\#\# 3. 几何学的解釈

方程式 \( x^2 = -3 \) に遭遇し、これは正または負の実数で表すことができない解を持ちませんでした。その後、方程式 \( x^2 = -1 \) を解釈できるなら、つまり \(\sqrt{(-1)} × \sqrt{(-1)} = -1\) となるような \(\sqrt{(-1)}\) を定義できれば、すべての困難が消えます。

今、特別な順序対 \((0, 0)\)、\((1, 0)\)、および \((0, 1)\) を考えましょう。

既に証明済みですが：
\[
(x, y) + (0, 0) = (x, y).
\]

さらに現在は：
\[
(x, y) × (0, 0) = (0, 0).
\]

したがって、加算と乗算の両方において、順序対 \((0, 0)\) は初等的な算術と代数におけるゼロの役割を果たします。上記の式を \( x + 0 = x \)、および \( x × 0 = 0 \) と比較してください。

また、\((1, 0)\) も考慮します：これは初等的な算術と代数における \(1\) の役割を果たします。これらの初等的科学では、\(1\) の特別な特性はすべての値の \(x\) に対して \( x × 1 = x \) であることです。現在の乗算法則によれば：
\[
(x, y) × (1, 0) = \{(x - 0), (y + 0)\} = (x, y).
\]

したがって、\((1, 0)\) は単位順序対です。

\#\#\#\# 4. 複素数の定義と特性

Finally consider \((0, 1)\): このものは \(\sqrt{(-1)}\) の意味を表します。したがって、\(\sqrt{(-1)} × \sqrt{(-1)} = -1\) という特性を持つべきです。現在、順序対の乗法法則により：
\[
(0, 1) × (0, 1) = \{(0 - 1), (0 + 0)\} = (-1, 0).
\]

\((1, 0)\) は単位対であり、\((-1, 0)\) は負の単位対です。したがって、\((0, 1)\) には望まれる特性があります。しかし、\(-1\) の平方根は \(±\sqrt{(-1)}\) と2つあります。次に \((0, -1)\) を考慮します；ここで再度、\((-1)^2 = 1\) を念頭に置いて：
\[
(0, -1) × (0, -1) = (-1, 0).
\]

したがって、\((0, -1)\) は \(\sqrt{(-1)}\) のもう一つの平方根です。従って、順序対 \((0, 1)\) と \((0, -1)\) は \(±\sqrt{(-1)}\) の順序対による解釈となります。

しかし、どちらがどの符号に対応するのでしょうか？\((0, 1)\) が \(+\sqrt{(-1)}\) に対応し、\((0, -1)\) が \(-\sqrt{(-1)}\) に対応するのか、それとも \((0, 1)\) が \(-\sqrt{(-1)}\) に対応し、\((0, -1)\) が \(+\sqrt{(-1)}\) に対応するのでしょうか？答えは、採用するシンボリズムを選択することに完全な自由があるということです。

\#\#\#\# 5. 複素数の種類

順序対は三種類に分類できます：

- **(i) 両方の成分がゼロでない「複素虚数」型 \((x, y)\)；**
- **(ii) 第二成分がゼロの「実数」型 \((x, 0)\)；**
- **(iii) 第一成分がゼロの「純虚数」型 \((0, y)\)。**

これらの種類間の関係を考慮します。まず、「複素虚数」対 \((x, y)\) と「実数」対 \((a, 0)\) の積を考えます：
\[
(a, 0) × (x, y) = (ax, ay).
\]

これは単に順序対 \((x, y)\) の各成分を正または負の実数 \(a\) で乗算するだけです。

次に、「複素虚数」対 \((x, y)\) と「純虚数」対 \((0, b)\) の積を考えます：
\[
(0, b) × (x, y) = (-by, bx).
\]

ここでの効果はより複雑で、幾何学的解釈を考慮する前にさらに3つの特別なケースを挙げてから理解するのが最も適切です。

第三に、「実」の偶対 \((a, 0)\) を「虚」の偶対 \((0, b)\) で掛けると、
\[
(a, 0) × (0, b) = (0, ab).
\]

第四に、「実」の偶対 \((a, 0)\) と \((a', 0)\) を掛けると、
\[
(a, 0) × (a', 0) = (aa', 0).
\]

第五に、「虚」の偶対 \((0, b)\) と \((0, b')\) を掛けると、
\[
(0, b) × (0, b') = (-bb', 0).
\]

次に幾何学的な解釈に移ります。

\#\#\#\# 6. 几何学的解釈

これらの定義と法則は、複素数の世界を理解するための重要な基盤となります。順序対 \((x, y)\) の各成分 \(x\) と \(y\) を考えると、\(x\) は実軸上に位置し、\(y\) は虚軸上に位置します。つまり、複素数 \(z = x + yi\) は平面上の点 \((x, y)\) に対応します。

この視覚的な解釈は、複素数の乗算が単純な回転と拡大を表すことを示します。例えば、\((0, 1)\) を考えると、これは虚単位 \(i\) の順序対に対応し、\((0, -1)\) は \(-i\) の順序対に対応します。

以上で、複素数の乗法とその特性について詳しく説明しました。これらの概念を理解することで、我々は複雑な数学的な問題にもより深く向き合えるでしょう。

\chapter{Chapter IX: Coordinate Geometry}

\#\#\# Chapter IX: 坐標幾何

坐標幾何的手法と概念は、前章でも使用されてきましたが、現在はそれらをより精密に検討する時です。これにより他の到達した概念に対する理解も深まるでしょう。本章および次章では、正負実数という概念に戻り、最後二章で導入された虚数については無視します。

我々は常に、平面において二つの軸$XOX'$と$YOY'$を用いて、どのような点$P$も正のまたは負の数$x$および$y$という一対の数で位置を定めることができるという考えを使用してきました。ここで（参照図13）$x$は長さ$OM$であり、$y$は長さ$PM$です。この概念は見かけ上単純ですが、座標幾何学の主要なアイデアです。その発見は数学的思想の歴史上重要な時代を示すものです。これは（既に述べた通り）哲学者デカルトによって生み出されました\index{Descartes}{}、彼はベッドで横になっている朝にこの重要な数学者の方法を見つけることに到達しました。

哲学者は数学的な知識が深ければ深いほど、その科学に最良のアイデアをもたらすことが多いです。一方で、それらの哲学者のうちほとんど例外なく、僅かなまたは急いで得られた数学に関する言及は全く価値がないものであり、 trivia か誤りであることが多くあります。事実、数学における最終的な概念は非常に単純で、子供でも理解できるほど簡単であり、哲学的な思考の範囲内に留まっています。おそらくその極めて単純さが誤りを生む原因かもしれません。我々はこのような単純な抽象的概念について考える習慣がなく、既成の思考の道からそれた瞬間には長年の訓練が必要でさえあります。

座標幾何の発見と、ほぼ同時期にProjective幾何学も発見したことは、知識の歴史において繰り返し確認されている事実を示しています。つまり、最も広く知られているトピックの中から、いくつかの最伟大的な発見がなされることがしばしばあります。

17世紀になると、幾何学はすでに紀元前2000年以上にわたって研究されていました。ギリシャ人による幾何学の隆盛期を始めるのが適切であれば、その起源はユーリドクスにまでさかのぼります。アレクサンドリア大学で教鞭を執ったユーリドクスは紀元前330年頃生まれで、彼は長い系列の前人作業を系統化し、発展させました。その一部は天才的な人物でした。

その後、世代を超えて数学者たちはこの分野の改善に努め続けました。しかし、その進歩を妨げる致命的な障害は、その研究が狭い集団の類似した出身と視点を持つ人々によって制限されていたことではありませんでした。むしろ、17世紀にはエジプト人やギリシャ人、アラブ人やドイツ人の頭脳を通じて幾何学が広く普及していました。

それでも、そのような多様な頭脳によって長い年月にわたって研究が行われた後でも、最も重要な秘密はまだ発見されていませんでした。

誰も初等幾何学の要素を学んだことなく、導き方となる方法の欠如を感じたことはないだろう。各命題は新たな独創性の展示によって証明される必要がある。そして、この特徴を持つ科学は、科学的な思考にとって重要な要素である方法性を欠いている。

座標幾何学の特別な点は、それが初めて方法を導入したという点にある。数学的科学の遠く離れた推論は、理論上重要なものではない。科学が完全なものになるまでには、その本質は、任意の範囲内のどのトピックについても容易に情報を得るための大規模で関連する方法の展示であるべきだ。

科学の成長は主に量ではなく、アイデアによるものであり、アイデアが増えるほど、書き出す価値のある推論は少なくなる。不幸にも、数学は教科書によって数多くの補助的な命題が繰り返し記述され、それらの重要性はより一般的な真理の特殊な例として吸収されて失われている。そして、既に強調したように、一般性は数学の魂である。

また、座標幾何学はもう一つの特徴を示す。それはすでに指摘された通り、数学的科学が発展するにつれて、互いに嵌め合わ合い、共通のアイデアを持つことだ。言い換えれば、数学の各分野は永続的な一般化プロセスを経験し、その過程で融合していく。この理由は、科学の本質からくるものであり、それが普遍性を持っていること、つまりすべての存在するものの性質に関連する一般的な真理を扱うことに由来する。

座標幾何学の関心点は、空間の科学として始まった幾何学と、数の科学として起源を持つ代数学を結びつける点にある。

現在、二つの科学の主なアイデアを思い返し、デカルトの座標法によってどのように関連付けるかを見ていきましょう。まず代数について考えてみましょう。虚数については触れず、正または負の符号を持つ実数のみを考えます。基本的なアイデアは、任意の数である変数数というものです。これは文字で表され、具体的な数値では表現されません。次に、変数間の関係性を考慮します。例えば、$x$ と $y$ が二つの変数の場合、これらは方程式 $x + y = 1$ または $x - y = 1$ によって関連付けられることも考えられます。あるいは無数の他の方法でも関連付けられることがあります。

これにより、代数形の概念を適用することができます。実際には、何か興味深い型の関係性を考えます。つまり、変数数という最初の概念から、数値の変数関係性という次に進んだ概念へと昇華します。したがって、$x + y = 1$ の関係を一般化して $ax + by = c$ の関係にします。ここで $a$、$b$ そして $c$ は文字であり、任意の数を表し、実際にはこれらも変数です。しかし、これらの変数が変数関係性を決定するものです。関係性が確定されると、それは変数数 $x$ と $y$ の間に関連付けられます。

$a$、$b$、そして $c$ などのような、関係性を決定するために使用される変数は「定数」またはパラメータと呼ばれます。この文脈での「定数」という用語の使用が最初に見えて不自然に思えるかもしれませんが、実際には非常に自然です。なぜなら、数学的な調査は $a$、$b$、$c$ が既に決定された後の $x$ と $y$ の関係性についての関係を調べることに関与しているからです。したがって、$x$ および $y$ 相対的に、「定数」である $a$、$b$、そして $c$ は実質的に定数となります。したがって、$ax + by = c$ は特定の代数形の一般的な例を表し、あるクラスに属する変数関係性を表します。

また、$x^{2} + y^{2} = 1$ を一般化して $ax^{2} + by^{2} = c$ とします。
さらに、$ax^{2} + 2hxy + by^{2} = c$ とします。
そして、さらに $ax^{2} + 2hxy + by^{2} + 2gx + 2fy = c$ とします。

またここでも、様々な代数的形式によって示される変動的な関係が導き出されます。

次に幾何学へと移りましょう。その科学の名前は、すぐに私たちの心を図形や図表、三角形や長方形、正方形や円などの考えに連想させます。これらの図形の単純な性質から始めて、より複雑な形状へと進化します。

\#\#\# 坐標幾何の基本

坐標幾何は、平面内の点を数値で表現する方法を提供します。例えば、点$P$が$x$軸上の長さ$OM$と$y$軸上の長さ$PM$を持つとき、点$P$の位置は$(x, y)$という一対の数で定義されます。この概念は非常に単純ですが、それが坐標幾何学の基礎となっています。

デカルトがこの発見をした背景には、彼がベッドで横になっている朝に突然閃いたという逸話があります。これは、哲学者として深い思考を行う中で、数学的な洞察を得た瞬間を象徴しています。哲学と数学は密接に関連しており、哲学的な思考が数学の発展にも大きく貢献します。

\#\#\# 代数と幾何学の融合

坐標幾何学は、代数と幾何学を結びつける重要な橋渡しとなりました。代数では、変数$x$や$y$を使って方程式を作ります。例えば、$x + y = 1$という方程式は、直線上の点$(x, y)$が満たす条件を表します。

さらに一般化すると、$ax + by = c$という形で方程式を定義できます。ここで$a$、$b$、そして$c$は定数（パラメータ）であり、これらは具体的な数値を取ることができます。これらの定数が変数$x$と$y$の関係性を決定します。

例えば、$x^2 + y^2 = 1$という方程式は円を表します。これを一般化すると、$ax^2 + by^2 = c$という形になります。さらに、$ax^2 + 2hxy + by^2 + 2gx + 2fy = c$という形も考えられます。

これらの式は、異なる形状の図形（円、楕円、双曲線など）を表すことができます。つまり、代数的な方程式が幾何学的な図形を生成するという原理が坐標幾何学の核心です。

\#\#\# まとめ

座標幾何学は、数学と哲学の融合点であり、新しい視点から既存の知識を再解釈します。デカルトの発見は、代数と幾何学を結びつける重要な橋渡しとなりました。この方法論は、数学的な問題解決に新たな可能性を開くとともに、科学全体の発展にも大きな影響を与えました。

座標幾何学を通じて、我々は空間と数の関係性をより深く理解することができます。これは、私たちが日常生活で直面する多くの問題を解くための強力なツールとなり得ます。

\chapter{Chapter X: Conic Sections}

\#\#\# Chapter X: Conic Sections

ギリシャの幾何学者たちは、直線と円で構成される図形の性質を深く探求した後、他の曲線にも焦点を移しました。彼らは放物線、双曲線、楕円などの放物面で切り取られる円錐表面の曲線に注目し始めました。この学科の発明者として信頼されるべき人物はメネラウスです（紀元前375年生～紀元前325年没）。彼はプラトンの弟子であり、アレクサンドロス大王の教師の一人でもありました。また、彼の他の教師の一人は哲学者アリストテレスでした。

メネラウスはアレクサンドロスに証明を短くするようにと頼まれたことがあります。しかし、メネラウスは次のように返答しました：「国には秘密の道や王様の道があるが、幾何学にはすべての人々のために一つだけの道しかない」。この回答はアレクサンドロスがすぐに理解した意味においては十分に真実でした。しかし、メネラウスが証明を短くすることは不可能だと考えていたら、彼は深刻な誤りを犯していたのです。そして、多くの現代の数学者は、ギリシャ人が放物線の性質を証明する方法を学ぶことを強いられたらひどく退屈することでしょう。

科学における概念の導入によって得られる力の増加を示す例としては、証明が短くなる過程と同時に概念の豊かさが成長する過程を観察することでよく理解できます。ある種の数学者は、主題のアイデアに耽るよりも早く重要な問題の証明に進むことを常に少し不満に感じることがあります。しかし、科学の歴史は彼に対して完全に対立しています。科学には王様の道がありますが、最初にそれを歩いたのは天才であり、国王ではありません。

円錐曲線が数学的に最初に現れた方法は次の通りでした：円錐（図\[fig.]{15}\参照）を考え、その頂点を$V$とし、底面が円形の$STU$上に立っています。例えば、電球に対して設置されたコンパクトな影はこのような表面の例です。今、$V$を通る「生成」する直線が全て延長されると、二つの円錐が得られます。そして$PQR$は$V$と反対側にある別の円形の交差点となります。円錐$CVC'$の軸はこれらの円の中心を通り、それぞれの平面に垂直であり、これらは互いに平行です。図において紙の平面の奥に存在する曲線の部分は点線で示され、平面またはその前方にある部分は継続線で示されています。

今、この二つの円錐が軸$CVC'$に対して垂直ではない平面上で切られるとします。あるいは、必ずしも垂直である必要はありません。その結果、以下の3つの状況が生じます：

1. 平面が一つの半円錐上に完全に存在する閉じた楕円形の曲線、例えば$ABA'B'$のように切れる場合があります。この場合、平面はもう片方の半円錐と全く交わらないことがあります。このような曲線は楕円と呼ばれ、それは楕円形の曲線です。円錐のこのような断面の特別な例として、平面が軸$CVC'$に対して垂直である場合、断面、例えば$STU$や$PQR$は円となります。したがって、円は楕円の特別な例と言えます。

2. 平面は、円錐の生成線の一つに沿って接平面を引くことができます。例えば図において、曲線 $D_{1}A_{1}D_{1}'$ の平面は生成線 $VS$ に沿って接する接平面と平行であり、この曲線は依然として半円錐の一つに制限されるが、閉じた卵形の曲線ではなくなり、生成線が頂点から延長されると無限に続く。このような円錐截痕をパラボラという。

3. 平面が二つの半円錐を両方とも切る場合、完全な曲線は二つの離れた部分、または「枝」と呼ばれる二つの部分からなります。図において、二つの枝 $G_{2}A_{2}G_{2}'$ と $L_{2}A_{2}'L_{2}'$ が組み合わさって曲線を形成しています。どちらの枝も閉じていないことがわかり、半円錐が頂点から延長されるにつれて無限に広がります。このような円錐截痕を双曲線という。

従って、三種類の円錐截痕があります：楕円、パラボラ、および双曲線。パラボラは、一定の条件を満たすより特殊な種であり、楕円と双曲線の間の極限的な場合であることが容易に理解できます。アポロニウス・ペルガエウス（紀元前260年頃生没）によって書かれた系統的な円錐截痕に関する論文が存在し、16世紀まで標準的な著作として受け入れられました。

これらの曲線の性質を調査することは、ギリシャの幾何学者にとっていかに不快で困難な作業であったかがすぐに明らかとなる。これらの曲線は平面曲線であるが、その調査には立体図形を透视的に描くことが含まれている。したがって、上記の図においても補助線をほとんど引いていないにもかかわらず、図は十分に複雑になっている。これらの曲線は平面曲線であり、これらを定義するためには平面内だけで済ませることができそうに思える。同時に、「立体」での定義では一貫した方法論が一つ存在し、それは円錐を平面で切るという方法で、楕円、双曲線、そして放物線の三種類を生み出す。同様に、「平面」での定義でも一貫した方法論が三種類に分かれることであるべきだ。これらの曲線は図形 \FigNum{16}, \FigNum{17}, および \FigNum{18} の曲線の形状と一致している。図形において点 $A$ と $A'$ は頂点と呼ばれ、線分 $AA'$ は主軸と呼ばれる。

パラボラ（\cf \Fig[fig.]{17}）には一つだけ頂点があることが注意されるべきである。アポロニウスは、楕円と双曲線（図形 \FigNum{16} および \FigNum{18}）において $PM^{2}$ に対する $AM·MA'$ の比が一定であり、放物線（\cf \Fig[fig.]{17}）においては $PM^{2}$ に対する $AM$ の比が一定であることを証明した。彼はこの事実に基づいて多くの仕事を進めている。明らかに我々は平面内だけで定義する一貫性のある方法論へと進んでいるが、まだ完全には到達していない。

\#\#\# 翻訳

円錐曲線が数学的に最初に現れた方法は次の通りでした：円錐（図\[fig.]{15}\参照）を考え、その頂点を$V$とし、底面が円形の$STU$上に立っています。例えば、電球に対して設置されたコンパクトな影はこのような表面の例です。今、$V$を通る「生成」する直線が全て延長されると、二つの円錐が得られます。そして$PQR$は$V$と反対側にある別の円形の交差点となります。円錐$CVC'$の軸はこれらの円の中心を通り、それぞれの平面に垂直であり、これらは互いに平行です。図において紙の平面の奥に存在する曲線の部分は点線で示され、平面またはその前方にある部分は継続線で示されています。

今、この二つの円錐が軸$CVC'$に対して垂直ではない平面上で切られるとします。あるいは、必ずしも垂直である必要はありません。その結果、以下の3つの状況が生じます：

1. 平面が一つの半円錐上に完全に存在する閉じた楕円形の曲線、例えば$ABA'B'$のように切れる場合があります。この場合、平面はもう片方の半円錐と全く交わらないことがあります。このような曲線は楕円と呼ばれ、それは楕円形の曲線です。円錐のこのような断面の特別な例として、平面が軸$CVC'$に対して垂直である場合、断面、例えば$STU$や$PQR$は円となります。したがって、円は楕円の特別な例と言えます。

2. 平面は、円錐の生成線の一つに沿って接平面を引くことができます。例えば図において、曲線 $D_{1}A_{1}D_{1}'$ の平面は生成線 $VS$ に沿って接する接平面と平行であり、この曲線は依然として半円錐の一つに制限されるが、閉じた卵形の曲線ではなくなり、生成線が頂点から延長されると無限に続く。このような円錐截痕をパラボラという。

3. 平面が二つの半円錐を両方とも切る場合、完全な曲線は二つの離れた部分、または「枝」と呼ばれる二つの部分からなります。図において、二つの枝 $G_{2}A_{2}G_{2}'$ と $L_{2}A_{2}'L_{2}'$ が組み合わさって曲線を形成しています。どちらの枝も閉じていないことがわかり、半円錐が頂点から延長されるにつれて無限に広がります。このような円錐截痕を双曲線という。

従って、三種類の円錐截痕があります：楕円、パラボラ、および双曲線。パラボラは、一定の条件を満たすより特殊な種であり、楕円と双曲線の間の極限的な場合であることが容易に理解できます。アポロニウス・ペルガエウス（紀元前260年頃生没）によって書かれた系統的な円錐截痕に関する論文が存在し、16世紀まで標準的な著作として受け入れられました。

これらの曲線の性質を調査することは、ギリシャの幾何学者にとっていかに不快で困難な作業であったかがすぐに明らかとなる。これらの曲線は平面曲線であるが、その調査には立体図形を透视的に描くことが含まれている。したがって、上記の図においても補助線をほとんど引いていないにもかかわらず、図は十分に複雑になっている。これらの曲線は平面曲線であり、これらを定義するためには平面内だけで済ませることができそうに思える。同時に、「立体」での定義では一貫した方法論が一つ存在し、それは円錐を平面で切るという方法で、楕円、双曲線、そして放物線の三種類を生み出す。同様に、「平面」での定義でも一貫した方法論が三種類に分かれることであるべきだ。これらの曲線は図形 \FigNum{16}, \FigNum{17}, および \FigNum{18} の曲線の形状と一致している。図形において点 $A$ と $A'$ は頂点と呼ばれ、線分 $AA'$ は主軸と呼ばれる。

パラボラ（\cf \Fig[fig.]{17}）には一つだけ頂点があることが注意されるべきである。アポロニウスは、楕円と双曲線（図形 \FigNum{16} および \FigNum{18}）において $PM^{2}$ に対する $AM·MA'$ の比が一定であり、放物線（\cf \Fig[fig.]{17}）においては $PM^{2}$ に対する $AM$ の比が一定であることを証明した。彼はこの事実に基づいて多くの仕事を進めている。明らかに我々は平面内だけで定義する一貫性のある方法論へと進んでいるが、まだ完全には到達していない。

\#\#\# 翻訳

図 \FigNum{16} および \FigNum{18} では、二点 $S$ と $S'$ を見ることができます。また、図 \Fig[diagram]{17} では一点 $S$ を見ることができます。これらは曲線の焦点であり、最も重要な点です。アポロニウスは、楕円の場合、点 $P$

\chapter{Chapter XI: Functions}

\#\#\# Chapter XI: Functions

数学における「関数」という言葉は、私たちの日常生活でも頻繁に使われています。例えば、「彼の気分は消化状態による関数である」という表現は、この数学的な意味で正確です。これは、彼の気分がどのように消化状態によって決まっているかを知るルールがあることを意味します。つまり、「関数」の概念は単純で、変数の値に対してどのように適用されるかを見るだけです。

まずいくつか具体的な例を考えましょう：ある列車が時速20マイルで走行している場合、任意の時間 \( t \)（時間）後の距離 \( s \)（マイル）は \( s = 20 × t \) で与えられます。そして \( s \) は \( t \) の関数と呼ばれ、\( 20 × t \) も同じ関数です。また、ジョンがトーマスより1歳年上である場合、トーマスの年齢 \( x \)（年）が任意のとき、ジョンの年齢 \( y \)（年）は \( y = x + 1 \) で与えられます。そして \( y \) も関数 \( x + 1 \) と同一です。

これらの例では、\( t \) と \( x \) を関数に現れる「引数」と呼びます。したがって \( t \) は関数 \( 20 × t \) の引数であり、\( x \) は関数 \( x + 1 \) の引数です。もし \( s = 20 × t \) で、\( y = x + 1 \) である場合、\( s \) と \( y \) はそれぞれ関数 \( 20 × t \) と関数 \( x + 1 \) の「値」と呼ばれます。

現在、一般の場合について述べるが、数学における関数を変動する二つの数の間の相関として定義できる。それぞれを「関数の引数」と「関数値」呼ばれるものとするとき、関数の引数にどのような値を代入しても、関数値は一意に（すなわち独特に）決定される。逆は必ずしも成り立たない。つまり、関数値が決定されても、その引数が一意に決定されるわけではない。

引数 \( x \) の他の関数として \( y = x^2 \), \( y = 2x^2 + 3x + 1 \), \( y = x \), \( y = \log x \), \( y = \sin x \）がある。このグループの最後二つの関数は、少しの代数学と三角法を理解している人にはすぐに識別できるだろう。その説明に時間を費やす価値はなく、例示のためだけに引用されているものである。

これまでまで、関数の一般的な定義は述べましたが、特定の関数の一連のみを挙げてきました。しかし、数学は一般的な方法で進行するため、任意の関数の一般的な概念を表すために符号化します。これは、\( x \) の関数として任意の関数を表すために \( F(x) \), \( f(x) \), \( g(x) \), \( \phi(x) \）などと記述することで達成されます。ここで、引数 \( x \) は括弧に入れられ、\( F \), \( f \), \( g \), \( \phi \) などのような文字が括弧の前に置かれることで関数を表します。

この表記法には欠点があります。つまり、これは明らかに単一の文字が変動する数字を表す慣習と衝突します；ここでは、括弧の前に置かれた \( F \), \( f \), \( g \), \( \phi \) などが変動する関数を表しています。意味が明確でない例を挙げると、直感や文脈に頼らなければなりません。混乱を避けるための一法は、関数のためにギリシャ文字（例えば上の \( \phi \））を使用することです；また、関数の文字として \( f \) と \( F \) （関数の頭文字）を使用し、他の変動する関数が必要な場合は隣接する文字（例：\( g \)）を用いることです。

これらの説明と注意事項に基づき、\( y = f(x) \) と書きます。これは \( y \) が引数 \( x \) の未定義の関数の値であることを表します；ここで \( f(x) \) は \( x + 1 \), \( x^2 - 2x + 1 \), \( \sin x \), \( \log x \），または単に \( x \) 自身を表すことができます。重要な点は、\( x \) が与えられたときには \( y \) もそれに従って確定的に定まるということです。この考えの一般性について十分に理解することが重要です。したがって \( y = f(x) \) の場合、例えば \( f(x) \) を \( x \) が整数であるとき \( f(x) = 0 \) で、それ以外の値を持つとき \( f(x) = 1 \）とすることができます。この意味での \( f(x) \) を用いて \( y = f(x) \) と置くと、\( x \) の値が整数であるときは \( y = 0 \) であり、それ以外の場合は \( y = 1 \) となります。したがって \( f(1) = 0 \), \( f(2) = 0 \), \( f(\frac{2}{3}) = 1 \), \( f(\sqrt{2}) = 1 \）などです。この \( f(x) \) の意味の選択は、引数 \( x \) の関数として一般的な定義に従って完全に良い関数を提供します。

結局のところ関数は、二つの変数間の一種の相関関係であるため、他の相関関係同様座標幾何学的なグラフで表されます。例えば第II章のFig. 2 は、引数 \( v \) と関数の値 \( p \) の関数 \( \frac{1}{v} \）のグラフです。この場合、グラフは正の \( v \) のみを描いており、その物理的な応用に意味を持つ唯一の値です。また第IX章のFig. 14 では、線分 \( AB \) 全体（両方向に無限延長）が引数 \( x \) と関数の値 \( y \) の関数 \( x + 1 \）のグラフであり、同じ図において無限な線分 \( A_1B \) は関数 \( 1 - x \）のグラフであり、線分 \( LOL' \) は引数 \( x \) と関数の値 \( y \) の関数 \( x \）のグラフです。

これらの関数は単純な代数的公式で表され、グラフによって表現できる。しかし、一部の関数については、詳細な説明なしではそのグラフによる表現は非常に誤解を招くものであり、甚至不可能である場合もある。例えば、引数 \( x \) が整数以外のすべての値に対して値が1となるが、\( x = 0, x = 1, x = 2,\ldots \）のような整数では値が0となる関数を考える。この関数のグラフ上の表現は、\( XOX' \)軸から1単位離れた平行線である直線 \( ABA' \)に相当する。しかし、\( x \) の値が \( 0, 1, 2, 3, 4,\ldots \）に対応する点 \( B, C_1, C_2, C_3, C_4,\ldots \)は省略され、代わりに軸 \( OX \) 上の点 \( O, B_1, B_2, B_3, B_4,\ldots \)が取り入れられる。グラフによる表現が不便であるだけでなく不可能な関数も存在し、そのような関数は高次の数学において重要であり、ここではそれらについてさらに詳しく触れることはない。

関数の最も重要な分類は連続関数と不連続関数との区別にあり\index{連続関数|EtSeq}\index{不連続関数|EtSeq}。引数が徐々に変化するときその値も徐々に変化する関数を連続関数、突然の跳躍によって値が変わる関数を不連続関数と呼ぶ。図\[fig.]{14}\ の第9章で示したグラフである \( x + 1 \) と \( 1 - x \）は連続関数であり、第2章で示された \( \frac{1}{v} \) も正の \( v \) の値を考える限りでは連続関数である。しかし、この章の図\[fig.]{20}\ に示した関数は不連続関数であり、その引数 \( x = 1, x = 2,\ldots \）

\chapter{Chapter XII: Periodicity in Nature}

\#\#\# Chapter XII: 周期性と自然

自然の全生涯は、周期的な出来事の存在によって支配されています。すなわち、類似性が高い互換的に繰り返される一連の出来事が存在することです。地球の回転は、連続する日を生み出します。どの程度抽象的に「日」と定義しても、各日は前の日に異なるものであることは否定できません。しかし、「日」の意味を十分に抽象化すれば、二つの日の性質的な違いは実用的な興味から遠ざかり、それぞれの日は地球の1回の回転の現象の再現として考えることができます。

また、地球が太陽周りを巡る経路は年間の季節の再現をもたらし、自然全体のすべての操作に別の周期性を課します。月相によるもう一つの基本的でない周期性もあります。現代の文明社会では、人工的な光によってこれらの月相は微々たるものですが、古代時代、日が灼熱し晴れ渡った気候の中で人間の生活は月明かりの存在に大きく影響されたようです。そのため、週や月による分類とその宗教的関連性はシリアや美索不達ミアからヨーロッパ諸民族へと広まりました。しかし、多くの国々では、月相に基づく独立した儀礼が見られます。

しかしながら、月の周期性は主に潮汐を通じて地球の歴史に大きな影響を与え、その明暗による周期性よりも重要な役割を果たしました。

私たちの身体的生活は本質的に周期的である。心臓の拍動と呼吸の再現がその主導を握っている。周期性の前提は、実際上生命自体の概念に根本的なものである。我々は進展する出来事に対して、「これ以前にも起こっていた」と言える状況を想像できない。経験が行動指針となるためには、過去の歴史と類似した新しい状況に常に直面することもなければ、時間量を測定する手段もなくなる。ある出来事が他のよりも前後関係にあるという単純な認識を超えて、私たちは次のように言えることができる：
三つの出来事 $A$,~$B$,~$C$ がこの順序で起こったので、$A$ は $B$ より前に、そして $B$ は $C$ より前に起こった。さらに、$A$ の発生から $B$ の発生までの時間の長さが、$B$ の発生から $C$ の発生までの時間の長さの二倍であることも言える。

現在の時間量は、自然な再現が介入した数を観察することに本質的に依存している。私たちは、$A$ と $B$ の間の時間を、何日、何ヶ月、または何年など、我々が参照したいタイプの再現によって、そのような数だけと言える。

実際には、文明が始まった段階では、これらの三つの時間測定方法は本当に異なるものだった。科学は文明や半文明国家にとって最初の重要な任務の一つであり、これらを一貫した計測に統合することである。この任務の全貌を理解することが必要だ。単にどの数日（例えば $365.25$ など）が某一年に相当するかだけでなく、さらに前もって次の年々にも同じ数日の長さが適用されることを決定しなければならない。

周期性が存在する世界を想像できるが、その周期性は互いに整合的ではない場合もある。ある年の間に $200$ 日で、他の年の間には $350$ 日であるなどということがあり得る。より重要な周期性の広範な一貫性の決定は自然科学の最初の段階だった。この一貫性は抽象的な直感的法則ではなく、経験によって保証された自然の観察事実によるものだ。

\#\#\# 周期性と時間測定

私たちが周期性を理解する際、時間の測定方法が異なる場合、同じ物体に対して異なる速度変動を得ることになります。例えば、1時間は1日の中の24分の1と定義し、列車が2時間にわたって平均速度で20マイル毎時で走行するという状況を考えてみましょう。次に、非常に不一致な時間測定法を考え、最初の1時間が次の1時間の2倍長であるとする場合があります。すると、この別の時間尺度では、列車の走行時間が二つの部分に分かれ、各部分で同じ距離、すなわち20マイルを走行しますが、最初の部分の持続時間は次の部分の持続時間の2倍です。したがって、列車の速度は一貫して均一ではなく、平均的には第二部分での速度は第一部分での速度の2倍となります。

地球上の日常生活におけるすべての通常の目的において、様々な天文的再発現は絶対的に一致していると見なすことができます。さらにこれらの一致性を仮定し、それによって物体の速度や速度変化も仮定すると、上記で考慮した運動の法則がほとんど正確に確認されます。しかし、いくつかの天文的な現象に関しては、完全には正確ではありません。しかし、惑星と恒星の回転および運動のわずかに異なる速度を仮定することで、法則が正確に確認されることを見つけることができます。

\#\#\# 周期性と共振

共振は動力学法則であり、すべての物体が自己の振動は特有の一定の時間で行われることがある。つまり、小さな振幅での単振り子はその形状や重心分布、長さによって特有の一定の時間を常に振動する。より複雑な物体は多くの振動方法を持つかもしれないが、それぞれの振動モードには独自の「周期」がある。これらの体の振動の周期は「自由」周期と呼ばれる。したがって単振り子は一つの周期しか持たない一方で、吊り橋は多数

\#\#\# 吊り橋の共振

吊り橋は多くの振動モードを持つ可能性があります。例えば、吊り橋は風に揺れたり、歩行者が通ったりすることで振動します。これらの振動はそれぞれ異なる「周期」を持ちます。しかし、特定の条件下では、一つの振動が他の振動と共振する可能性があります。共振とは、二つの関連する状況が同じ周期を持つときに生じる現象です。

例えば、吊り橋が風に揺れると、その振動は吊り橋の固有の「自由」周期に合わせます。しかし、歩行者が吊り橋を歩くとき、歩行者の足拍子が吊り橋の自由周期と一致すると、共振が生じます。この場合、吊り橋の振動は増幅され、結果として大きな揺れが発生します。

\#\#\# 周期性と時間尺度

したがって、精密な天文観測によって検出されるその不一致はどう対処すればよいのか？明らかに私たちは、これらの現象のいずれかが等しい時間を表すという随意的な仮定に至る。例えば、すべての日が等長であるとするか、またはすべての年が等長であるとするか。

しかし、そうではない：何かしらの仮定が必要だが、天文学者が時間の尺度を決定する全過程において基礎となる仮定は、運動の法則が完全に遵守されているという点にある。つまり、地球の回転や太陽系の他の天体の運動が一定の規則性を持つと仮定します。

この方法がどのように機能するかを説明する前に、そのような上位の一貫性が人間の体の特性的な再現と一致していた場合、時計の調整は医師たちに委ねられたことだろうという興味深い観察を行うのは面白い。実際、心臓の拍動や呼吸のリズムなど、生物に関連した再現は他の再現と比較して急速に変化する。一方で、大きな安定性を有し、互いに非常に正確に一致する重要な再現は、地球全体の運動や類似の天体の運動によるものである。

\#\#\# まとめ

周期性は自然の基本的な特徴であり、私たちの生活にも深く関わっています。時間の測定方法が異なる場合でも、一定の法則に基づいて宇宙と地球上の物事が動いていることがわかります。共振や自由周期など、これらの概念を理解することで、より深い洞察を得ることができます。科学はこれらの法則を明らかにし、私たちの生活をより正確に把握する手助けをしてくれます。

以上が本章の解説講義です。ご質問や補足があればお知らせください。

\end{document}
